\documentclass[11pt]{article}
\usepackage[a4paper,margin=1in]{geometry}
\usepackage{amsmath,amssymb,amsthm}
\usepackage{mathtools}      % \coloneqq
\usepackage{bm}
\usepackage{enumitem}
\usepackage{algorithm}
\usepackage{algpseudocode}
\usepackage[colorlinks=true,linkcolor=blue,citecolor=blue]{hyperref}
\usepackage[capitalize,noabbrev]{cleveref}

% ---- Theorem environments: independent counters, per section ----
\theoremstyle{plain}
\newtheorem{theorem}{Theorem}[section]
\newtheorem{lemma}{Lemma}[section]
\newtheorem{proposition}{Proposition}[section]
\newtheorem{corollary}{Corollary}[section]
\theoremstyle{definition}
\newtheorem{definition}{Definition}[section]
\newtheorem{assumption}{Assumption}[section]
\theoremstyle{remark}
\newtheorem{remark}{Remark}[section]

\crefname{assumption}{Assumption}{Assumptions}
\Crefname{assumption}{Assumption}{Assumptions}

% ---- Macros ----
\DeclarePairedDelimiter{\norm}{\lVert}{\rVert}
\DeclarePairedDelimiter{\abs}{\lvert}{\rvert}
\newcommand{\R}{\mathbb{R}}
\newcommand{\N}{\mathbb{N}}
\newcommand{\E}{\mathbb{E}}
\newcommand{\Prob}{\mathbb{P}}
\newcommand{\Unif}{\mathrm{Unif}}
\newcommand{\Vol}{\mathrm{Vol}}
\renewcommand{\d}{\mathrm{d}}
\DeclareMathOperator{\tr}{tr}
\DeclareMathOperator{\diag}{diag}
\DeclareMathOperator{\blockdiag}{blockdiag}
\newcommand{\lmin}{\lambda_{\min}}
\newcommand{\grad}{\nabla}
\newcommand{\Hess}{\nabla^2}
\newcommand{\inner}[2]{\langle #1,\, #2\rangle}
\newcommand{\eps}{\varepsilon}
\newcommand{\epsg}{\eps_{g}}
\newcommand{\epsH}{\eps_{H}}
\newcommand{\Bmin}{B_{\min}}
\newcommand{\mup}{\mu_{+}}
\newcommand{\mum}{\mu_{-}}
\newcommand{\vp}{v_{+}}
\newcommand{\vm}{v_{-}}
\newcommand{\gp}{\gamma_{+}}
\newcommand{\gm}{\gamma_{-}}
\newcommand{\Mmat}{\mathbf{M}}
\newcommand{\xt}{\tilde{x}}
\newcommand{\xs}{x^{\star}}
\newcommand{\Rstar}{R_{0}^{\star}}
\newcommand{\Otil}{\widetilde{O}}
\newcommand{\sg}{\sigma_{\mathrm{err}}}

\allowdisplaybreaks

% ============================================================
\title{\bf A Perturbed Constant-Momentum Method\\
Escapes Strict Saddle Points}

\author{%
  Benjamin Séguin\\
  LAMSADE - Miles\\%
}
\date{\today}
% ============================================================

\begin{document}
\maketitle

\begin{abstract}
We study a perturbed momentum method with a constant momentum parameter,
a fixed step size, and a restart safeguard, and prove that it converges
to an approximate second-order stationary point of a smooth nonconvex
objective. Writing
$L_{1}$ and $L_{2}$ for the gradient- and Hessian-Lipschitz constants
and $(\epsg,\epsH)$ for the first- and second-order tolerances, the
method finds a point with $\norm{\grad f(x)}\le\epsg$ and
$\lmin(\Hess f(x))\ge-\epsH$ in
\[
  O\!\left(\frac{L_{1}\,\Delta_{f}}{\epsg^{2}}\right)
  \;+\;
  \Otil\!\left(\frac{L_{1}^{3}L_{2}^{2}\,\Delta_{f}}{\epsH^{4}}\right)
\]
iterations with high probability, where $\Delta_{f}$ bounds the initial
suboptimality and $\Otil$ hides polylogarithmic factors in the
dimension and the failure probability. Under the standard scaling
$\epsH=\sqrt{L_{2}\epsg}$ this matches the $\Otil(\epsg^{-2})$ rate of
perturbed gradient descent. Two ingredients drive the analysis. First,
a restart test on the momentum direction forces every
search direction to be a descent direction with a quantitative angle
bound, which yields a per-step decrease and an improve-or-localize
estimate. Second, around a strict saddle we linearise the
\emph{difference} of two coupled runs that share their perturbation but
are offset along the escape direction; the difference obeys a
non-symmetric matrix recursion whose dominant eigenvalue we identify
explicitly, and a thin-slab volume argument shows that the set of
perturbations from which the method fails to escape is negligible.
\end{abstract}

% ============================================================
\section{Introduction}
\label{sec:intro}

Nonconvex optimisation problems are pervasive in modern machine
learning, and although locating a global minimiser is intractable in
the worst case, a more modest target has proved both useful and
attainable: convergence to a stationary point. A point with small
gradient may, however, be a saddle point or even a local maximum, and
the loss surfaces arising in applications are often densely populated
by such points. It is therefore natural to ask for convergence to a
\emph{second-order} stationary point, where in addition to a small
gradient the Hessian is approximately positive semidefinite, ruling out
the strict saddle points at which $\lmin(\Hess f)<0$.

For gradient descent (GD), the analysis of Jin, Ge, Netrapalli, Kakade
and Jordan~\cite{pmlr-v70-jin17a} shows that a single random perturbation,
injected whenever the gradient is small, suffices to escape strict
saddles, and that the perturbed method reaches an approximate
second-order stationary point in $\Otil(\eps^{-2})$ iterations---the
same rate, up to logarithmic factors, at which GD reaches a first-order
stationary point. The companion work of Jin, Netrapalli and
Jordan~\cite{pmlr-v75-jin18a} extends the perturbation idea to Nesterov's
accelerated gradient method and isolates two reusable tools: an
\emph{improve-or-localize} principle that confines slowly-progressing
iterates to a small ball, and a \emph{coupling} argument that bounds the
volume of the ``stuck'' region by tracking two trajectories separated
along the escape direction.

The present paper carries this programme to a method with
\emph{constant} momentum. We study the iteration
$x_{k+1}=x_{k}+\alpha\,d_{k}$, $d_{k+1}=-g_{k+1}+\beta\,d_{k}$ with a
fixed coefficient $\beta$---equivalently, the heavy-ball recursion
$x_{k+1}=x_{k}-\alpha\grad f(x_{k})+\beta(x_{k}-x_{k-1})$~\cite{POLYAK19641}---together
with a restart safeguard that keeps the momentum direction aligned with
the negative gradient, and a perturbation that is injected when the
gradient is small. Two features distinguish the analysis from the
gradient-descent case. The restart keeps every step a descent step
despite the momentum; and, unlike gradient descent, whose linearisation
is symmetric, the constant-momentum iteration linearises to a
\emph{non-symmetric} propagator that does not diagonalise in the
eigenbasis of the Hessian, which we control by an explicit spectral
analysis of the two-trajectory difference. Keeping $\beta$ constant is a
deliberate restriction: it isolates the geometric core of the argument
and sets aside the gradient-dependent coefficient of genuine nonlinear
conjugate gradient, whose denominator $\norm{\grad f}^{2}$ vanishes at
stationary points (see \Cref{sec:conclusion}).

Our analysis rests on two ideas. The first is local: a restart test of
the form $\norm{g_{k}+d_{k}}\le\nu\norm{g_{k}}$ makes every search
direction a descent direction with a quantitative angle bound, valid
both on momentum and on restart steps;
together with the fixed step size $\alpha=1/L_{1}$ this gives a per-step
decrease proportional to $\norm{g_{k}}^{2}$ and an improve-or-localize
estimate. The second is the coupling: near a strict saddle we follow two
runs that share the perturbation but are offset along the smallest
curvature direction, linearise their difference, and show that the
non-symmetric propagator has a dominant real eigenvalue
$\mup\ge1+\alpha\epsH>1$ that drives the two runs apart, while a
thin-slab volume argument bounds the set of perturbations from which the
method stays stuck. \Cref{sec:setup} fixes the setting, the method and
the main theorem; \Cref{sec:descent} establishes the descent and
localization estimates; \Cref{sec:recurrence} sets up the linearized
two-trajectory recursion; \Cref{sec:spectral} analyses the propagator
and identifies the escape eigenvalue; \Cref{sec:induction} shows the
geometric growth survives the nonlinearity; \Cref{sec:escape} turns this
into a high-probability escape guarantee; and \Cref{sec:complexity}
assembles the iteration complexity.

We emphasise that the method is not accelerated: the rate we obtain
matches perturbed gradient descent~\cite{pmlr-v70-jin17a}, not the
$\Otil(\eps^{-7/4})$ of perturbed AGD~\cite{pmlr-v75-jin18a}. Our contribution is
to show that the perturbation-and-coupling machinery extends to a
constant-momentum iteration, whose non-symmetric propagator we analyse
explicitly, once the momentum direction is safeguarded by restarts.

\paragraph{Related work.}
% Convergence of first-order methods to second-order stationary points has
% a substantial literature. Cubic regularisation~\cite{NesterovPolyak2006}
% and trust-region methods reach an $\eps$-second-order point in
% $O(\eps^{-3/2})$ iterations but require Hessian access. Among
% gradient-based methods, Ge, Huang, Jin and Yuan~\cite{Ge2015} gave the
% first polynomial guarantee for noisy gradient descent, sharpened by Jin
% et al.~\cite{Jin2017} to $\Otil(\eps^{-2})$ with only polylogarithmic
% dimension dependence. Hessian-vector-product and nested-loop schemes
% attain $\Otil(\eps^{-7/4})$~\cite{CarmonDuchi2018,RoyerWright2018}, a
% rate also achieved by the single-loop perturbed AGD
% of~\cite{JNJ2018}. The stable-manifold viewpoint of Lee et
% al.~\cite{Lee2019} shows that many first-order methods avoid saddles
% from almost every initialisation, but without an explicit rate. Momentum methods themselves are classical---the heavy-ball method dates
% to Polyak~\cite{Polyak1964}---and their behaviour near saddle points has
% been studied in the quadratic and accelerated settings. A fixed-step
% nonlinear conjugate gradient method with first-order complexity
% guarantees was analysed by Chan-Renous-Legoubin and
% Royer~\cite{ChanRoyer2022}; extending the present second-order analysis
% to that gradient-dependent coefficient is our motivating goal
% (\Cref{sec:conclusion}). Our two-trajectory argument is the coupling
% mechanism of~\cite{Jin2017,JNJ2018}; the non-symmetry of the
% constant-momentum propagator is what forces the spectral treatment of
% \Cref{sec:spectral}.

\paragraph{Notation.}
Throughout, $\norm{\cdot}$ is the Euclidean norm for vectors and the
spectral norm for matrices, and $\lmin(\cdot)$ is the smallest
eigenvalue of a symmetric matrix. We write $g_{k}\coloneqq\grad
f(x_{k})$, $\Hess f$ for the Hessian, and $B_{0}(r)$ for the Euclidean
ball of radius $r$ centred at the origin. We use $O(\cdot)$ and
$\Theta(\cdot)$ in the usual way and $\Otil(\cdot)$ to hide
polylogarithmic factors in all problem parameters. The symbol $\xt$
denotes the snapshot at which a perturbation is added (the candidate
saddle), and $\xs$ a global minimiser. Constants introduced along the
way are collected in \Cref{tab:ledger}.

% ============================================================
\section{Setting, Algorithm and Main Result}
\label{sec:setup}

We consider the unconstrained problem $\min_{x\in\R^{d}}f(x)$ under the
following standing assumptions.

\begin{assumption}\label{ass:smooth}
$f\colon\R^{d}\to\R$ is twice continuously differentiable and
$L_{1}$-gradient-Lipschitz: $\norm{\grad f(x)-\grad f(y)}\le
L_{1}\norm{x-y}$ for all $x,y$.
\end{assumption}

\begin{assumption}\label{ass:hesslip}
$f$ is $L_{2}$-Hessian-Lipschitz: $\norm{\Hess f(x)-\Hess f(y)}\le
L_{2}\norm{x-y}$ for all $x,y$.
\end{assumption}

\begin{assumption}\label{ass:bounded}
$f$ is bounded below, and $\Delta_{f}\coloneqq f(x_{0})-\inf_{x}f(x)<\infty$.
\end{assumption}

\begin{definition}[Approximate second-order stationary point]
\label{def:sosp}
Given tolerances $\epsg,\epsH>0$, a point $x$ is an
$(\epsg,\epsH)$-second-order stationary point ($2$OSP) of $f$ if
\[
  \norm{\grad f(x)}\le\epsg
  \qquad\text{and}\qquad
  \lmin\!\big(\Hess f(x)\big)\ge-\epsH .
\]
\end{definition}

The single-tolerance formulation of~\cite{jipmlr-v70-jin17a,pmlr-v75-jin18a} sets
$\epsH=\sqrt{L_{2}\,\eps}$ and $\epsg=\eps$; we keep the two tolerances
separate to follow the formulation of \cite{RoyerChan2022}.

\paragraph{The method.}
The method is the constant-momentum iteration
\[
  x_{k+1}=x_{k}+\alpha\,d_{k},
  \qquad
  d_{k}=-g_{k}+\beta\,d_{k-1},
\]
with a \emph{fixed} momentum coefficient $\beta\in(0,1]$ and step size
$\alpha=1/L_{1}$, together with two safeguards. Eliminating the search
direction gives the equivalent heavy-ball
form~\cite{POLYAK19641}
\[
  x_{k+1}=x_{k}-\alpha\,\grad f(x_{k})+\beta\,(x_{k}-x_{k-1}).
\]
\begin{itemize}[leftmargin=*,itemsep=2pt]
\item \emph{Restart test.} The momentum direction is accepted
only if it is sufficiently aligned with the negative gradient,
\begin{equation}\label{eq:restart-test}
  \norm{g_{k}+d_{k}}\le\nu\,\norm{g_{k}},
\end{equation}
with slack $\nu\in(0,\sqrt{5}-2)$; otherwise the method restarts,
setting $d_{k}=-g_{k}$. We write $\mathcal{N}$ for the steps on which the
momentum direction is accepted (momentum steps) and $\mathcal{R}$ for the
restart steps; on $\mathcal{R}$ the test~\eqref{eq:restart-test} holds
with $\nu=0$.
\item \emph{Perturbation.} When $\norm{g_{k}}\le\epsg$ and no
perturbation has been added in the last $T$ iterations, the iterate is
perturbed, $x_{k}\leftarrow x_{k}+\xi$ with $\xi\sim\Unif(B_{0}(R))$, and
the direction is reset.
\end{itemize}
The step size, restart slack, momentum $\beta$, perturbation radius $R$,
window length $T$ and gradient trigger $\epsg$ are fixed in
\Cref{sec:complexity}; their dependence on the problem data is recorded
in \Cref{tab:ledger}.

\begin{algorithm}[t]
\caption{PR-CMOM: Perturbed Restarted Constant Momentum
$\mathrm{PR\text{-}CMOM}(x_0,\,L_1,L_2,\,\varepsilon_g,\varepsilon_H,\,\beta,\,\nu,\,\delta_{\mathrm{tot}},\,\Delta_f)$}
\label{alg:prncg}
\begin{algorithmic}[1]
\State \textbf{Derived constants} (resolve the constraint system of Section~\ref{sec:complexity}):
\Statex \quad $\alpha \gets 1/L_1$;\quad
        $\Lambda \gets \log\!\frac{d\, L_1 L_2 \Delta_f}{\varepsilon_g\,\varepsilon_H\,\delta_{\mathrm{tot}}}$;\quad
        $c(\nu)\gets (1-\nu)-\tfrac{(1+\nu)^2}{2}$
\Statex \quad $\mu_+ \gets$ larger root of $\mu^2-(1+\beta+\alpha\varepsilon_H)\mu+\beta$;\quad
        resolve $(T,S,F,r)$ and per-episode $\delta=\delta_{\mathrm{tot}}/\bar K$
\State $d_0 \gets -\nabla f(x_0)$;\quad $t_{\mathrm{noise}} \gets -T-1$
\For{$t = 0,1,2,\dots$}
    \If{$\|\nabla f(x_t)\| \le \varepsilon_g$ \textbf{and} $t - t_{\mathrm{noise}} > T$}
        \Comment{gradient small: test for negative curvature}
        \State $\tilde{x} \gets x_t$;\quad $t_{\mathrm{noise}} \gets t$
        \State $x_t \gets \tilde{x} + \xi_t$,\quad $\xi_t \sim \mathrm{Unif}\big(\mathcal{B}^d(0,r)\big)$
               \Comment{perturb}
        \State $d_t \gets -\nabla f(x_t)$
               \Comment{reset momentum at the snapshot}
    \EndIf
    \If{$t - t_{\mathrm{noise}} = T$ \textbf{and} $f(x_t) - f(\tilde{x}) > -F$}
        \State \Return $\tilde{x}$
               \Comment{no decrease over the window: $\tilde{x}$ is an $(\varepsilon_g,\varepsilon_H)$-2OSP}
    \EndIf
    \State $x_{t+1} \gets x_t + \alpha\, d_t$
           \Comment{fixed step $\alpha=1/L_1$}
    \State $g_{t+1} \gets \nabla f(x_{t+1})$
           \Comment{\textbf{constant} momentum}
    \State $\hat{d}_{t+1} \gets -g_{t+1} + \beta\, d_t$
           \Comment{candidate conjugate direction}
    \If{$\|\hat{d}_{t+1} + g_{t+1}\| \le \nu\,\|g_{t+1}\|$}
        \State $d_{t+1} \gets \hat{d}_{t+1}$
               \Comment{keep momentum (restart test passes)}
    \Else
        \State $d_{t+1} \gets -g_{t+1}$
               \Comment{restart: drop momentum}
    \EndIf
\EndFor
\end{algorithmic}
\end{algorithm}

Our main result is the following; its proof occupies
\Cref{sec:descent,sec:recurrence,sec:spectral,sec:induction,sec:escape}
and is assembled in \Cref{sec:complexity}.

\begin{theorem}[Main result]
\label{thm:main}
Let \Cref{ass:smooth,ass:hesslip,ass:bounded} hold and fix a total
failure probability $\delta_{\mathrm{tot}}\in(0,1)$. There is a choice of
parameters (\Cref{sec:complexity}) under which, provided each escape window is a momentum window
(\Cref{def:window}), the method visits
an $(\epsg,\epsH)$-second-order stationary point with probability at
least $1-\delta_{\mathrm{tot}}$ after at most
\[
  N \;=\;
  \underbrace{O\!\left(\frac{L_{1}\,\Delta_{f}}{C(\nu)\,\epsg^{2}}\right)}_{
      \text{large-gradient phase}}
  \;+\;
  \underbrace{\Otil\!\left(\frac{L_{1}^{3}L_{2}^{2}\,\Delta_{f}}{\epsH^{4}}\right)}_{
      \text{escape phase}}
  \;=\;
  O\!\big(\epsg^{-2}\big)+\Otil\!\big(\epsH^{-4}\big)
\]
iterations, where $C(\nu)=\tfrac{1}{2}(1-4\nu-\nu^{2})>0$ and $\Otil$
hides factors polylogarithmic in $d$, $1/\delta_{\mathrm{tot}}$ and the
problem constants.
\end{theorem}

\begin{remark}
Under the coupling $\epsH=\sqrt{L_{2}\,\epsg}$ of \Cref{def:sosp},
$\epsH^{-4}=(L_{2}\epsg)^{-2}$ and the escape phase costs
$\Otil(\epsg^{-2})$, so the overall rate is $\Otil(\epsg^{-2})$,
matching perturbed gradient descent~\cite[Theorem 7]{pmlr-v70-jin17a}.
\end{remark}

\begin{remark}
The escape-phase bound is established on momentum windows
(\Cref{def:window}); whether restarts can be ruled out inside an escape
window is discussed in \Cref{sec:conclusion}.
\end{remark}

\paragraph{Proof strategy near a saddle.}
The large-gradient phase is handled by a counting argument
(\Cref{sec:complexity}). The crux is the escape phase. Suppose the
method perturbs at a snapshot $\xt$ with
$\lmin(\Hess f(\xt))\le-\epsH$. We run two coupled copies from
$x_{0}=\xt+\xi$ and $x_{0}'=x_{0}+R_{0}v_{\min}$, sharing the
perturbation but offset by $R_{0}$ along the smallest curvature
direction $v_{\min}$. \emph{(I)} Their difference obeys a linear
recursion whose homogeneous part grows like $\mup^{k}R_{0}$ with
$\mup\ge1+\alpha\epsH>1$, and the nonlinear residual stays a fixed
fraction of it, so the gap grows geometrically. \emph{(II)} A copy that
fails to escape within $T$ steps is confined, by improve-or-localize, to
a ball of radius $S$; two confined copies cannot be farther than
$2S+R_{0}$ apart, which caps the separation of any two stuck starts at a
threshold $\Rstar$. \emph{(III)} The stuck set is therefore a slab of
width $\Rstar$ along $v_{\min}$, and \emph{(IV)} a uniform perturbation
lands in it with probability at most
$(\Rstar/R)\sqrt{(d+1)/(2\pi)}$, which we drive below $\delta$.
Repeating over the at most $\Otil(\epsH^{-3})$ saddle episodes and
counting the large-gradient steps gives \Cref{thm:main}.

% ============================================================
\section{A Descent Direction and Localization}
\label{sec:descent}

We first turn the restart test into the two estimates that the rest of
the analysis rests on: a quantitative descent property of the search
direction, and the consequent per-step decrease.

\begin{lemma}[Direction bounds]
\label{lem:dirbounds}
On every iteration, whether a momentum step ($\mathcal{N}$) or a
restart step ($\mathcal{R}$),
\begin{equation}\label{eq:dirbounds}
  \norm{d_{k}}\le(1+\nu)\norm{g_{k}}
  \qquad\text{and}\qquad
  \inner{d_{k}}{g_{k}}\le-(1-\nu)\norm{g_{k}}^{2}.
\end{equation}
\end{lemma}

\begin{proof}
The accepted direction satisfies $\norm{g_{k}+d_{k}}\le\nu\norm{g_{k}}$
(on restart steps with $\nu=0$, since $d_{k}=-g_{k}$). For the norm,
\[
  \norm{d_{k}}=\norm{(d_{k}+g_{k})-g_{k}}
  \le\norm{d_{k}+g_{k}}+\norm{g_{k}}
  \le(1+\nu)\norm{g_{k}}.
\]
For the inner product, write $d_{k}=(d_{k}+g_{k})-g_{k}$, so that
\[
  \inner{d_{k}}{g_{k}}
  =-\norm{g_{k}}^{2}+\inner{d_{k}+g_{k}}{g_{k}}
  \le-\norm{g_{k}}^{2}+\norm{d_{k}+g_{k}}\,\norm{g_{k}}
  \le(\nu-1)\norm{g_{k}}^{2}. \qedhere
\]
\end{proof}

\begin{lemma}[Per-step decrease]
\label{lem:descent}
Let \Cref{ass:smooth} hold and let $\alpha=1/L_{1}$. For every
non-perturbation step,
\begin{equation}\label{eq:descent}
  f(x_{k+1})\le f(x_{k})-\frac{C(\nu)}{L_{1}}\,\norm{g_{k}}^{2},
  \qquad
  C(\nu)=(1-\nu)-\frac{(1+\nu)^{2}}{2}=\frac{1-4\nu-\nu^{2}}{2},
\end{equation}
and $C(\nu)>0$ if and only if $\nu<\sqrt{5}-2$.
\end{lemma}

\begin{proof}
By $L_{1}$-smoothness and $x_{k+1}=x_{k}+\alpha d_{k}$,
\[
  f(x_{k+1})\le f(x_{k})+\alpha\,\inner{g_{k}}{d_{k}}
  +\frac{L_{1}}{2}\alpha^{2}\norm{d_{k}}^{2}.
\]
Insert the two bounds of \Cref{lem:dirbounds}:
\[
  f(x_{k+1})\le f(x_{k})
  -\alpha(1-\nu)\norm{g_{k}}^{2}
  +\frac{L_{1}}{2}\alpha^{2}(1+\nu)^{2}\norm{g_{k}}^{2}
  = f(x_{k})+\norm{g_{k}}^{2}\!\left[-\alpha(1-\nu)
  +\frac{L_{1}}{2}\alpha^{2}(1+\nu)^{2}\right].
\]
With $\alpha=1/L_{1}$ the bracket equals
$-\tfrac{1}{L_{1}}\big[(1-\nu)-\tfrac{1}{2}(1+\nu)^{2}\big]
=-C(\nu)/L_{1}$. Finally
$C(\nu)=\tfrac12(2-2\nu-(1+\nu)^{2})=\tfrac12(1-4\nu-\nu^{2})$ is
positive exactly when $\nu^{2}+4\nu-1<0$, i.e.\ $\nu<\sqrt5-2$.
\end{proof}

The decrease~\eqref{eq:descent} also controls how far the iterates can
travel, which is the improve-or-localize principle: either the objective
drops appreciably, or the iterates stay in a small ball.

\begin{lemma}[Improve or localize]
\label{lem:iol}
Under the setting of \Cref{lem:descent}, suppose only
momentum and restart steps are taken on $[0,T]$. Then for every
$k\le T$,
\begin{equation}\label{eq:iol}
  \norm{x_{k}-x_{0}}\le(1+\nu)
  \sqrt{\frac{T\,\big(f(x_{0})-f(x_{T})\big)}{L_{1}\,C(\nu)}}.
\end{equation}
In particular, if $f(x_{0})-f(x_{T})\le F$ then
$\norm{x_{k}-x_{0}}\le S$ for all $k\le T$, where
\begin{equation}\label{eq:Sdef}
  S\;\coloneqq\;(1+\nu)\sqrt{\frac{T F}{L_{1}\,C(\nu)}}.
\end{equation}
\end{lemma}

\begin{proof}
Using $x_{k+1}-x_{k}=\alpha d_{k}$, the bound
$\norm{d_{k}}\le(1+\nu)\norm{g_{k}}$, and~\eqref{eq:descent} in the
form $\norm{g_{k}}^{2}\le\frac{L_{1}}{C(\nu)}(f(x_{k})-f(x_{k+1}))$,
\[
  \norm{x_{k+1}-x_{k}}^{2}
  =\alpha^{2}\norm{d_{k}}^{2}
  \le\alpha^{2}(1+\nu)^{2}\norm{g_{k}}^{2}
  \le\frac{(1+\nu)^{2}}{L_{1}C(\nu)}\big(f(x_{k})-f(x_{k+1})\big),
\]
since $\alpha^{2}=1/L_{1}^{2}$. Summing telescopes:
\[
  \sum_{j=0}^{k-1}\norm{x_{j+1}-x_{j}}^{2}
  \le\frac{(1+\nu)^{2}}{L_{1}C(\nu)}\big(f(x_{0})-f(x_{k})\big).
\]
By the triangle and Cauchy--Schwarz inequalities,
\[
  \norm{x_{k}-x_{0}}
  \le\sum_{j=0}^{k-1}\norm{x_{j+1}-x_{j}}
  \le\sqrt{k}\left(\sum_{j=0}^{k-1}\norm{x_{j+1}-x_{j}}^{2}\right)^{1/2}
  \le(1+\nu)\sqrt{\frac{k\,(f(x_{0})-f(x_{k}))}{L_{1}C(\nu)}},
\]
and $k\le T$, $f(x_{0})-f(x_{k})\le f(x_{0})-f(x_{T})\le F$ give the
claim.
\end{proof}

% ============================================================
\section{Linearizing Two Coupled Trajectories}
\label{sec:recurrence}
 
We now set up the coupling. Fix a snapshot $\xt$ at which a perturbation
is added, and consider two runs of the method,
\[
  x_{0}=\xt+\xi,\qquad x_{0}'=x_{0}+R_{0}\,v_{\min},
\]
that share the perturbation $\xi$ but are offset by $R_{0}$ along the
minimal curvature direction $v_{\min}$ of $H\coloneqq\Hess f(\xt)$. Write
$\Delta x_{k}=x_{k}-x_{k}'$ and $\Delta d_{k}=d_{k}-d_{k}'$, and study
\[
  w_{k}\coloneqq\begin{pmatrix}\Delta x_{k}\\[2pt]\Delta d_{k-1}\end{pmatrix}
  \in\R^{2d}.
\]
The key is to write the gradient as its linear part plus a remainder,
$g_{k}=Hx_{k}+e_{k}$ with $e_{k}\coloneqq\grad f(x_{k})-Hx_{k}$, and to
control the remainder through the Hessian-Lipschitz constant.

\begin{definition}[Momentum window]\label{def:window}
A block of $T$ iterations following a perturbation is a \emph{momentum
window} if, for both coupled trajectories, the restart
test~\eqref{eq:restart-test} fails to trigger at every step, so that
$d_{k}=-g_{k}+\beta\,d_{k-1}$ holds with the same fixed $\beta$
throughout. Equivalently, the propagator $\Mmat$ below is constant
across the window.
\end{definition}

\noindent The escape phase is analysed on momentum windows; this is the
analogue of the ``only accelerated steps, no negative-curvature
exploitation'' regime of~\cite{pmlr-v75-jin18a}. A restart replaces $\Mmat$ by
the matrix obtained by setting $\beta=0$, so the single-matrix growth
estimates of \Cref{sec:spectral} no longer apply; discharging
\Cref{def:window}---showing restarts are inactive near a saddle---is left
to future work (\Cref{sec:conclusion}).

\begin{lemma}[Difference recursion]
\label{lem:recurrence}
Along a momentum window (\Cref{def:window}), with $\alpha=1/L_{1}$, and the constant momentum $\beta$,
\begin{equation}\label{eq:recurrence}
  w_{k+1}=\Mmat\,w_{k}+\mathcal{R}_{k},
  \qquad
  \Mmat=\begin{pmatrix}I-\alpha H & \alpha\beta I\\[2pt] -H & \beta I\end{pmatrix},
  \qquad
  \mathcal{R}_{k}=\begin{pmatrix}-\alpha\,\Delta e_{k}\\[2pt]-\Delta e_{k}\end{pmatrix},
\end{equation}
where $\Delta e_{k}=\grad f(x_{k})-\grad f(x_{k}')-H\Delta x_{k}$.
Moreover, if all iterates of both runs stay within $S$ of their starts
and within $R$ of $\xt$ at the start, then
\begin{equation}\label{eq:residual}
  \norm{\mathcal{R}_{k}}\le\sg\,\norm{w_{k}},
  \qquad \sg\coloneqq \sqrt{1+\alpha^{2}}\;L_{2}(S+R).
\end{equation}
\end{lemma}

\begin{proof}
From $x_{k+1}=x_{k}+\alpha d_{k}$ and
$d_{k}=-g_{k}+\beta d_{k-1}=-(Hx_{k}+e_{k})+\beta d_{k-1}$ we obtain, on
differencing the two runs,
\[
  \Delta d_{k}=-H\Delta x_{k}-\Delta e_{k}+\beta\,\Delta d_{k-1},
  \qquad
  \Delta x_{k+1}=\Delta x_{k}+\alpha\,\Delta d_{k}.
\]
Substituting the first into the second,
\[
  \Delta x_{k+1}
  =(I-\alpha H)\Delta x_{k}+\alpha\beta\,\Delta d_{k-1}-\alpha\,\Delta e_{k},
\]
which is the top block of~\eqref{eq:recurrence}; the bottom block is the
displayed expression for $\Delta d_{k}$. This proves the recursion.
 
For the residual, set $E(x)\coloneqq\grad f(x)-Hx$, whose Jacobian is
$J_{E}(x)=\Hess f(x)-H$, so that by \Cref{ass:hesslip},
$\norm{J_{E}(x)}\le L_{2}\norm{x-\xt}$. Writing
$y_{t}=x_{k}'+t(x_{k}-x_{k}')$ and applying the fundamental theorem of
calculus,
\[
  \Delta e_{k}=E(x_{k})-E(x_{k}')
  =\Big(\int_{0}^{1}J_{E}(y_{t})\,\d t\Big)\Delta x_{k}.
\]
Hence
\[
  \norm{\Delta e_{k}}
  \le\norm{\Delta x_{k}}\int_{0}^{1}\norm{J_{E}(y_{t})}\,\d t
  \le L_{2}\norm{\Delta x_{k}}\int_{0}^{1}\norm{y_{t}-\xt}\,\d t
  \le L_{2}\,\max\{\norm{x_{k}-\xt},\norm{x_{k}'-\xt}\}\,\norm{\Delta x_{k}},
\]
using convexity of $t\mapsto\norm{y_{t}-\xt}$. Finally
$\norm{x_{k}-\xt}\le\norm{x_{k}-x_{0}}+\norm{x_{0}-\xt}\le S+R$, and the
same for the primed run, so
$\norm{\Delta e_{k}}\le L_{2}(S+R)\norm{\Delta x_{k}}$. The residual is
the stacked vector $\mathcal{R}_{k}=(-\alpha\Delta e_{k},-\Delta e_{k})$,
whose norm is $\norm{\mathcal{R}_{k}}=\sqrt{1+\alpha^{2}}\,\norm{\Delta e_{k}}$
exactly; combined with $\norm{\Delta x_{k}}\le\norm{w_{k}}$ this gives
$\norm{\mathcal{R}_{k}}\le\sqrt{1+\alpha^{2}}\,L_{2}(S+R)\norm{w_{k}}$,
which is~\eqref{eq:residual}.
\end{proof}

% ============================================================
\section{Spectral Analysis of the Propagator}
\label{sec:spectral}

The recursion~\eqref{eq:recurrence} is driven by the non-symmetric
matrix $\Mmat$. Because $\Mmat$ is a polynomial in $H$, it
block-diagonalises in the eigenbasis of $H$ into $2\times2$ blocks,
which we analyse one eigen-direction at a time.

\begin{lemma}[Block reduction]
\label{lem:block}
Let $H=VD V^{\top}$ with $V=[v_{1},\dots,v_{d}]$ and
$D=\diag(\lambda_{1},\dots,\lambda_{d})$. In the reordered basis
$(\Delta x^{(1)},\Delta d^{(1)},\dots,\Delta x^{(d)},\Delta d^{(d)})$,
the propagator is block-diagonal,
$\Mmat\cong\blockdiag(B_{1},\dots,B_{d})$, with
\begin{equation}\label{eq:blocks}
  B_{i}=\begin{pmatrix}1-\alpha\lambda_{i} & \alpha\beta\\[2pt]-\lambda_{i} & \beta\end{pmatrix},
  \qquad \det B_{i}=\beta,\qquad \tr B_{i}=1-\alpha\lambda_{i}+\beta.
\end{equation}
\end{lemma}

\begin{proof}
Conjugating $\Mmat$ by $\diag(V,V)$ replaces $H$ by $D$ in every
block of~\eqref{eq:recurrence}, giving
$\bigl(\begin{smallmatrix}I-\alpha D&\alpha\beta I\\-D&\beta
I\end{smallmatrix}\bigr)$. As all four sub-blocks are diagonal, grouping
the $i$-th coordinate of the $x$- and $d$-parts yields the $2\times2$
block $B_{i}$ in~\eqref{eq:blocks}; the determinant and trace are read
off directly.
\end{proof}

\begin{lemma}[Powers of a $2\times2$ block]
\label{lem:cayley}
Let $B$ be a $2\times2$ matrix with eigenvalues $\mu_{1},\mu_{2}$
(possibly equal or complex) and spectral radius
$\rho=\max\{\abs{\mu_{1}},\abs{\mu_{2}}\}$. For every $n\in\N$,
\[
  B^{n}=a_{n}B+b_{n}I,\qquad
  a_{n}=\frac{\mu_{1}^{n}-\mu_{2}^{n}}{\mu_{1}-\mu_{2}},\qquad
  b_{n}=-\mu_{1}\mu_{2}\,\frac{\mu_{1}^{n-1}-\mu_{2}^{n-1}}{\mu_{1}-\mu_{2}},
\]
and consequently
\begin{equation}\label{eq:Cbound}
  \norm{B^{n}}\le C(B)\,\rho^{n},
  \qquad
  C(B)\coloneqq\frac{2\,(\norm{B}+\rho)}{\abs{\mu_{1}-\mu_{2}}}.
\end{equation}
\end{lemma}

\begin{proof}
By the Cayley--Hamilton theorem $B^{2}=(\mu_{1}+\mu_{2})B-\mu_{1}\mu_{2}I$,
so $B^{n}=a_{n}B+b_{n}I$ for scalars $a_{n},b_{n}$. Evaluating on the two
eigenvalues, $\mu_{i}^{n}=a_{n}\mu_{i}+b_{n}$, and solving the $2\times2$
linear system gives the stated $a_{n},b_{n}$. Bounding
$\abs{\mu_{1}^{n}-\mu_{2}^{n}}\le2\rho^{n}$ and
$\abs{\mu_{1}\mu_{2}}\,\abs{\mu_{1}^{n-1}-\mu_{2}^{n-1}}
\le\abs{\mu_{1}}\,2\rho^{n}$ yields
$\abs{a_{n}}\le2\rho^{n}/\abs{\mu_{1}-\mu_{2}}$ and
$\abs{b_{n}}\le\abs{\mu_{1}}\,2\rho^{n}/\abs{\mu_{1}-\mu_{2}}$, hence
$\norm{B^{n}}\le\abs{a_{n}}\norm{B}+\abs{b_{n}}
\le\frac{2\rho^{n}}{\abs{\mu_{1}-\mu_{2}}}(\norm{B}+\rho)$.
\end{proof}

\subsection{The escape block}

The block with the most negative curvature is the engine of the escape.
We show it has a real eigenvalue strictly larger than one.

\begin{lemma}[Escape eigenvalue]
\label{lem:escapeeig}
Let $\lambda_{\min}\le-\epsH$ and let $\Bmin=B(\lambda_{\min})$ as
in~\eqref{eq:blocks}, with characteristic polynomial
$p(\mu)=\mu^{2}-t\mu+\beta$, where
$t=\tr\Bmin=1+\beta+\alpha\abs{\lambda_{\min}}$. Then $\Bmin$ has real
eigenvalues $\mum<1<\mup$ with
\begin{equation}\label{eq:mup-lb}
  \mup\;\ge\;1+\alpha\abs{\lambda_{\min}}\;\ge\;1+\alpha\epsH.
\end{equation}
\end{lemma}

\begin{proof}
Since $\lambda_{\min}<0$, $t=1+\beta+\alpha\abs{\lambda_{\min}}>1+\beta$.
Evaluating the characteristic polynomial at $1$,
\[
  p(1)=1-t+\beta=-\alpha\abs{\lambda_{\min}}<0,
\]
so $p$ has a root on each side of $1$, in particular $p$ has two real
roots and $\mum<1<\mup$. Writing $a=\alpha\abs{\lambda_{\min}}$ and
evaluating at $1+a$,
\[
  p(1+a)=(1+a)^{2}-(1+\beta+a)(1+a)+\beta
  =(1+a)\big[(1+a)-(1+\beta+a)\big]+\beta=-\beta(1+a)+\beta=-\beta a\le0.
\]
As $p(\mu)\to+\infty$ for $\mu\to\infty$ and $p(1+a)\le0$, the larger
root satisfies $\mup\ge1+a=1+\alpha\abs{\lambda_{\min}}\ge1+\alpha\epsH$.
\end{proof}

\begin{lemma}[No block grows faster]
\label{lem:nofaster}
Let $\beta>0$. For every eigenvalue
$\lambda\in[\lambda_{\min},L_{1}]$ of $H$, the spectral radius of
$B(\lambda)$ satisfies $\rho(\lambda)\le\mup$. Consequently, for the
propagator,
\begin{equation}\label{eq:Mbound}
  \norm{\Mmat^{n}}\le C_{0}\,\mup^{n},\qquad C_{0}\coloneqq\max_{i}C(B_{i}).
\end{equation}
\end{lemma}

\begin{proof}
By~\eqref{eq:blocks} the determinant $\det B(\lambda)=\beta$ is constant
while $t(\lambda)=1-\alpha\lambda+\beta$ is strictly decreasing in
$\lambda$, hence maximal at $\lambda=\lambda_{\min}$. The eigenvalues are
$\mu_{\pm}(\lambda)=\tfrac12\big(t(\lambda)\pm\sqrt{t(\lambda)^{2}-4\beta}\big)$.
If $t(\lambda)^{2}<4\beta$ the eigenvalues are complex conjugates with
$\abs{\mu}^{2}=\det=\beta$, so $\rho(\lambda)=\sqrt{\beta}$; we have
$\rho(\lambda)=\sqrt\beta = \sqrt{\mup \mum} \leq \mup$. If $t(\lambda)^{2}\ge4\beta$ the
larger root
$\rho(\lambda)=\tfrac12\big(t(\lambda)+\sqrt{t(\lambda)^{2}-4\beta}\big)$
is increasing in $t(\lambda)$, hence maximal at $\lambda=\lambda_{\min}$,
where it equals $\mup$. In both cases $\rho(\lambda)\le\mup$. Since
$\Mmat^{n}$ is block-diagonal,
$\norm{\Mmat^{n}}=\max_{i}\norm{B_{i}^{n}}\le\max_{i}C(B_{i})\rho_{i}^{n}
\le C_{0}\mup^{n}$ by~\eqref{eq:Cbound}.
\end{proof}

\subsection{Growth of the homogeneous part}

We decompose the difference state as $w_{k}=h_{k}+p_{k}$, the
homogeneous part $h_{k}=\Mmat^{k}w_{0}$ and the residual
$p_{k}=\sum_{\tau=0}^{k-1}\Mmat^{k-1-\tau}\mathcal{R}_{\tau}$, so that
$p_{0}=0$ and $p_{k+1}=\Mmat p_{k}+\mathcal{R}_{k}$. With $w_{0}$ chosen
in the escape plane, the homogeneous part is governed entirely by
$\Bmin$.

\begin{lemma}[Homogeneous growth]
\label{lem:homgrowth}
Let $w_{0}=(R_{0},0)$ lie in the two-dimensional escape plane and write,
in the eigenbasis of $\Bmin$, $v_{\mu}=(\beta-\mu,\lambda_{\min})$,
$\zeta=\mum/\mup\in[0,1)$. Fix $\theta\in(0,1)$ and let $k_{0}$ be the
least index with $\zeta^{k_{0}}\norm{\vm}\le\theta\norm{\vp}$. Then for
all $k\ge k_{0}$,
\begin{equation}\label{eq:hgrowth}
  \norm{h_{k}}\ge c_{1}\,\mup^{k}R_{0},
  \qquad c_{1}\coloneqq(1-\theta)\,\frac{\norm{\vp}}{\mup-\mum},
\end{equation}
and the position block alone satisfies
$\abs{\Delta x_{k}^{\mathrm{hom}}}\ge c_{1}'\,\mup^{k}R_{0}$ with
$c_{1}'=(1-\theta)\,\abs{\beta-\mup}/(\mup-\mum)$. Moreover
$\norm{h_{k}}\le c_{h}\,\mup^{k}R_{0}$ with $c_{h}=C(\Bmin)$.
\end{lemma}

\begin{proof}
The eigenvectors of $\Bmin$ solve $(\Bmin-\mu I)v=0$; from the second
row, $-\lambda_{\min}v^{(1)}+(\beta-\mu)v^{(2)}=0$, so
$v_{\mu}=(\beta-\mu,\lambda_{\min})$. Expand $w_{0}=\gp\vp+\gm\vm$. The
second coordinate gives $0=\gp\lambda_{\min}+\gm\lambda_{\min}$, hence
$\gm=-\gp$; the first gives
$R_{0}=\gp(\beta-\mup)+\gm(\beta-\mum)=\gp(\mum-\mup)$, so
$\abs{\gp}=R_{0}/(\mup-\mum)\neq0$. Thus $w_{0}$ has a nonzero component
along the unstable eigenvector. Therefore
\[
  h_{k}=\Bmin^{k}w_{0}
  =\gp\mup^{k}\vp+\gm\mum^{k}\vm
  =\gp\,\mup^{k}\big(\vp-\zeta^{k}\vm\big),
\]
and
\[
  \norm{h_{k}}
  \ge\abs{\gp}\,\mup^{k}\big(\norm{\vp}-\zeta^{k}\norm{\vm}\big)
  \ge(1-\theta)\,\abs{\gp}\,\norm{\vp}\,\mup^{k}
  =(1-\theta)\frac{\norm{\vp}}{\mup-\mum}\,\mup^{k}R_{0}
\]
for $k\ge k_{0}$. Projecting onto the position coordinate replaces
$\norm{\vp}$ by $\abs{\beta-\mup}$, giving the bound on
$\abs{\Delta x_{k}^{\mathrm{hom}}}$. The upper bound is immediate from
$\norm{h_{k}}=\norm{\Bmin^{k}w_{0}}\le C(\Bmin)\mup^{k}\norm{w_{0}}$
and $\norm{w_{0}}=R_{0}$.
\end{proof}

% ============================================================
\section{Geometric Growth of the Gap}
\label{sec:induction}

It remains to show that the residual $p_{k}$ does not cancel the growth
of $h_{k}$. The mechanism is an induction: assuming
$\norm{p_{\tau}}\le c\norm{h_{\tau}}$ up to step $k$, the bound on
$\norm{\mathcal{R}_{\tau}}$ keeps the accumulated residual a fraction of
$h_{k+1}$, provided $\sg=\sqrt{1+\alpha^{2}}\,L_{2}(S+R)$ is small enough, which is achievable because
$S$ and $R$ are at our disposal.

\begin{lemma}[Geometric growth survives]
\label{lem:induction}
Along a momentum window (\Cref{def:window}), suppose $c\in(0,1)$ and
\begin{equation}\label{eq:sigma-cond}
  \sg\;\le\;\frac{c\,c_{1}\,\mup}{C_{0}\,c_{h}\,(1+c)\,(T+1)}\;=:\;\frac{\Lambda}{T+1}.
\end{equation}
Then $\norm{p_{k}}\le c\,\norm{h_{k}}$ for all $k\le T$, and therefore
\begin{equation}\label{eq:gap-grows}
  \norm{w_{T}}\ge(1-c)\,c_{1}\,\mup^{T}R_{0},
  \qquad
  \norm{\Delta x_{T}}\ge\big(c_{1}'-c\,c_{h}\big)\,\mup^{T}R_{0}.
\end{equation}
\end{lemma}

\begin{proof}
We argue by induction on $k$. The base case $p_{0}=0$ is immediate.
Assume $\norm{p_{\tau}}\le c\norm{h_{\tau}}$ for all $\tau\le k$, so that
$\norm{w_{\tau}}\le(1+c)\norm{h_{\tau}}\le(1+c)c_{h}\mup^{\tau}R_{0}$ by
\Cref{lem:homgrowth}. Using
$p_{k+1}=\sum_{\tau=0}^{k}\Mmat^{k-\tau}\mathcal{R}_{\tau}$,
the propagator bound~\eqref{eq:Mbound}, the residual
bound~\eqref{eq:residual} and the inductive hypothesis,
\begin{align*}
  \norm{p_{k+1}}
  &\le\sum_{\tau=0}^{k}\norm{\Mmat^{k-\tau}}\,\norm{\mathcal{R}_{\tau}}
  \le\sum_{\tau=0}^{k}C_{0}\,\mup^{k-\tau}\,\sg\,\norm{w_{\tau}}\\
  &\le\sum_{\tau=0}^{k}C_{0}\,\mup^{k-\tau}\,\sg\,(1+c)\,c_{h}\,\mup^{\tau}R_{0}
  = C_{0}\,c_{h}\,\sg\,(1+c)\,R_{0}\,(k+1)\,\mup^{k}.
\end{align*}
By~\eqref{eq:sigma-cond} and $k+1\le T+1$,
$C_{0}c_{h}\sg(1+c)(k+1)\le c\,c_{1}\mup$, so
$\norm{p_{k+1}}\le c\,c_{1}\,\mup^{k+1}R_{0}\le c\,\norm{h_{k+1}}$,
the last step by the lower bound~\eqref{eq:hgrowth}. This closes the
induction. Finally
$\norm{w_{T}}\ge\norm{h_{T}}-\norm{p_{T}}\ge(1-c)\norm{h_{T}}
\ge(1-c)c_{1}\mup^{T}R_{0}$, and projecting onto the position block,
$\norm{\Delta x_{T}}\ge\abs{\Delta x_{T}^{\mathrm{hom}}}
-\norm{p_{T}}\ge c_{1}'\mup^{T}R_{0}-c\,c_{h}\mup^{T}R_{0}$.
\end{proof}

\begin{remark}
Positivity of the position growth coefficient requires
$c<c_{1}'/c_{h}$, which is built into the admissible range
$c\in(0,\min\{1,c_{1}'/c_{h}\})$ of \Cref{tab:ledger}.
\end{remark}

% ============================================================
\section{Escaping Strict Saddle Points}
\label{sec:escape}

We now combine the geometric growth of the gap with the confinement of
stuck trajectories. The decomposition of the perturbation along and
orthogonal to $v_{\min}$ is $\xi=\xi_{1}+\xi_{\perp}$ with
$\xi_{1}=\inner{\xi}{v_{\min}}\in\R$ and $\xi_{\perp}\in v_{\min}^{\perp}
\cong\R^{d-1}$. Define the localized and stuck sets
\[
  \mathcal{X}_{\mathrm{loc}}
  =\big\{\xi\in B_{0}(R):\norm{x_{k}-x_{0}}\le S\ \forall\,0\le k\le T\big\},
  \quad
  \mathcal{X}_{\mathrm{stuck}}
  =\big\{\xi\in B_{0}(R):f(x_{T})-f(x_{0})>-F\big\}.
\]
By improve-or-localize, a perturbation from which the method fails to
decrease the objective by $F$ is localized:
$\mathcal{X}_{\mathrm{stuck}}\subseteq\mathcal{X}_{\mathrm{loc}}$. It
therefore suffices to bound $\Prob(\xi\in\mathcal{X}_{\mathrm{loc}})$.

\begin{lemma}[Threshold separation]
\label{lem:threshold}
If two stuck trajectories start at $x_{0}$ and $x_{0}'=x_{0}+R_{0}v_{\min}$,
then $\norm{\Delta x_{T}}\le2S+R_{0}$. Combined with~\eqref{eq:gap-grows}
this forces
\begin{equation}\label{eq:Rstar}
  R_{0}\;\le\;\Rstar
  \;\coloneqq\;\frac{2S}{(c_{1}'-c\,c_{h})\,\mup^{T}-1}.
\end{equation}
\end{lemma}

\begin{proof}
If both runs are stuck, each is confined by \Cref{lem:iol} to a ball of
radius $S$ around its start, so
\[
  \norm{\Delta x_{T}}
  \le\norm{x_{T}-x_{0}}+\norm{x_{0}-x_{0}'}+\norm{x_{0}'-x_{T}'}
  \le S+R_{0}+S=2S+R_{0}.
\]
On the other hand~\eqref{eq:gap-grows} gives
$\norm{\Delta x_{T}}\ge(c_{1}'-c\,c_{h})\mup^{T}R_{0}$. Hence
$(c_{1}'-c\,c_{h})\mup^{T}R_{0}\le2S+R_{0}$, i.e.\
$R_{0}\big[(c_{1}'-c\,c_{h})\mup^{T}-1\big]\le2S$, which
is~\eqref{eq:Rstar}.
\end{proof}

\begin{lemma}[Escape with high probability]
\label{lem:escape}
Suppose $\norm{\grad f(\xt)}\le\epsg$ and
$\lmin(\Hess f(\xt))\le-\epsH$, and that a perturbation has not been
added in the previous $T$ steps, and that the next $T$ steps form a
momentum window (\Cref{def:window}). If the parameters satisfy
\begin{equation}\label{eq:Tcond}
  T\;\ge\;\frac{1}{\log\mup}\,
  \log\!\left[\frac{1}{c_{1}'-c\,c_{h}}
  \left(1+\frac{2S}{\delta R}\sqrt{\frac{d+1}{2\pi}}\right)\right]
\end{equation}
and $R$ is chosen so that $\epsg R+\tfrac{L_{1}}{2}R^{2}\le\eta F$ with
$\eta\le\tfrac12$, then with probability at least $1-\delta$,
\begin{equation}\label{eq:netdecrease}
  f(x_{T})-f(\xt)\le-(1-\eta)F\le-\tfrac12F.
\end{equation}
\end{lemma}

\begin{proof}
\emph{Volume of the stuck set.} By \Cref{lem:threshold}, any two stuck
perturbations lying on a line parallel to $v_{\min}$ are separated by at
most $\Rstar$, so the one-dimensional Lebesgue measure of
$\mathcal{X}_{\mathrm{loc}}$ along each such line is at most $\Rstar$. By
Fubini, splitting $x=(x^{(1)},x^{(\perp)})$ with $x^{(1)}$ along
$v_{\min}$,
\[
  \Vol(\mathcal{X}_{\mathrm{loc}})
  =\int_{B^{d-1}_{0}(R)}\!\!\Big(\int 
  \mathbf{1}\{\mathcal{X}_{\mathrm{loc}}\}\,\d x^{(1)}\Big)\d x^{(\perp)}
  \le\int_{B^{d-1}_{0}(R)}\!\!\Rstar\,\d x^{(\perp)}
  =\Rstar\,\Vol\!\big(B^{d-1}_{0}(R)\big).
\]
Using $\Vol(B^{n}_{0}(R))=\pi^{n/2}R^{n}/\Gamma(\tfrac n2+1)$ and
$\Gamma(\tfrac d2+1)/\Gamma(\tfrac d2+\tfrac12)\le\sqrt{(d+1)/2}$,
\[
  \Prob(\xi\in\mathcal{X}_{\mathrm{loc}})
  =\frac{\Vol(\mathcal{X}_{\mathrm{loc}})}{\Vol(B^{d}_{0}(R))}
  \le\frac{\Rstar}{R\sqrt\pi}\cdot\frac{\Gamma(\tfrac d2+1)}{\Gamma(\tfrac d2+\tfrac12)}
  \le\frac{\Rstar}{R}\sqrt{\frac{d+1}{2\pi}}.
\]
Requiring the right-hand side to be at most $\delta$ is, by the
definition~\eqref{eq:Rstar} of $\Rstar$, equivalent to
$\tfrac{2S}{\delta R}\sqrt{(d+1)/(2\pi)}+1\le(c_{1}'-c\,c_{h})\mup^{T}$,
i.e.\ to~\eqref{eq:Tcond}. Hence
$\Prob(\xi\in\mathcal{X}_{\mathrm{stuck}})
\le\Prob(\xi\in\mathcal{X}_{\mathrm{loc}})\le\delta$.

\emph{Net decrease.} On the complementary event, the trajectory leaves
the stuck set, so $f(x_{T})-f(x_{0})\le-F$. The perturbation itself
costs, by $L_{1}$-smoothness,
\[
  f(x_{0})-f(\xt)\le\inner{\grad f(\xt)}{\xi}+\frac{L_{1}}{2}\norm{\xi}^{2}
  \le\epsg R+\frac{L_{1}}{2}R^{2}\le\eta F.
\]
Adding the two, $f(x_{T})-f(\xt)\le-F+\eta F=-(1-\eta)F\le-\tfrac12F$
for $\eta\le\tfrac12$.
\end{proof}

% ============================================================
\section{Iteration Complexity}
\label{sec:complexity}
 
The escape lemma assumed a constraint system relating
$T,S,F,R,\delta$ and the growth constants. We record it, show it is
solvable, and then assemble the complexity bound and prove
\Cref{thm:main}.
 
\subsection{The constraint system}
 
Collecting the conditions from the previous sections, the operational
quantities must satisfy
\begin{align}
  &\text{(C1)}\quad S=(1+\nu)\sqrt{\tfrac{TF}{L_{1}C(\nu)}}
  \iff F=\frac{L_{1}C(\nu)}{(1+\nu)^{2}}\cdot\frac{S^{2}}{T},
  \label{eq:C1}\\
  &\text{(C2)}\quad \sg=\sqrt{1+\alpha^{2}}\,L_{2}(S+R)\le\frac{\Lambda}{T+1},
  \qquad \Lambda=\frac{c\,c_{1}\,\mup}{C_{0}\,c_{h}\,(1+c)},
  \label{eq:C2}\\
  &\text{(C3)}\quad L_{2}(S+R)\le\tfrac12\epsH,
  \label{eq:C3}\\
  &\text{(C4)}\quad \epsg R+\tfrac{L_{1}}{2}R^{2}\le\eta F,
  \label{eq:C4}\\
  &\text{(C5)}\quad T=\frac{1}{\log\mup}
  \log\!\Big[\tfrac{1}{c_{1}'-c\,c_{h}}\big(1+\tfrac{2S}{\delta R}
  \sqrt{\tfrac{d+1}{2\pi}}\big)\Big].
  \label{eq:C5}
\end{align}
Here $C(\nu),\Bmin,\mu_{\pm},c_{1},c_{1}',c_{h},C_{0}$ are the derived
quantities of \Cref{tab:ledger}, all functions of the problem data
$d,L_{1},L_{2},\Delta_{f}$ and the free parameters
$\epsg,\epsH,\beta,\nu,\theta,c,\eta,\delta$ (recall $\alpha=1/L_{1}$).
 
\begin{proposition}[Solvability]
\label{prop:solvable}
The system \eqref{eq:C1}--\eqref{eq:C5} admits a solution with
$R>0$ and $S>0$. Concretely, setting
$\Psi\coloneqq S+R=\Lambda/(\sqrt{1+\alpha^{2}}\,L_{2}(T+1))$
and
\[
  R=\min\!\left\{\frac{\eta F}{2\epsg},\ \sqrt{\frac{\eta F}{L_{1}}},\
  \frac{\Lambda}{2\sqrt{1+\alpha^{2}}\,L_{2}(T+1)}\right\},
  \qquad S=\Psi-R\ge\frac{\Lambda}{2\sqrt{1+\alpha^{2}}\,L_{2}(T+1)}>0,
\]
verifies (C2) with equality and (C4) by splitting each summand below
$\tfrac12\eta F$; (C3) follows from (C2) once $T\ge2\Lambda/\epsH-1$,
which holds because $\Lambda=\Theta(1)$ and $T=\Theta(L_{1}\phi/\epsH)\gg1/\epsH$;
and (C1) then determines $F$.
\end{proposition}
 
\begin{proof}
Set (C2) to equality, giving $\Psi=\Lambda/(\sqrt{1+\alpha^{2}}\,L_{2}(T+1))$.
The first two terms in the definition of $R$ enforce (C4): since
$\tfrac{L_{1}}{2}R^{2}\le\tfrac12\eta F$ when $R\le\sqrt{\eta F/L_{1}}$
and $\epsg R\le\tfrac12\eta F$ when $R\le\eta F/(2\epsg)$, their sum is
at most $\eta F$. The third term enforces $R\le\tfrac12\Psi$, when
$S=\Psi-R\ge\tfrac12\Psi>0$; any split fraction in $(0,1)$ would serve,
the requirement being only $R>0$ (a genuine perturbation) and $S>0$ (a
non-vacuous localization). With (C2) at equality,
$\sg=\sqrt{1+\alpha^{2}}\,L_{2}\Psi=\Lambda/(T+1)$, so
$L_{2}(S+R)=\Lambda/(\sqrt{1+\alpha^{2}}\,(T+1))\le\Lambda/(T+1)$ and (C3)
holds once $\Lambda/(T+1)\le\epsH/2$, i.e.\ $T+1\ge2\Lambda/\epsH$; as
$\Lambda=\Theta(1)$ and $T=\Theta(L_{1}\phi/\epsH)$ with $\phi\gg1$, this
holds. Finally (C1) defines $F$ from the chosen $S,T$.
\end{proof}
 
\begin{proposition}[Scales]
\label{prop:scales}
With the choice of \Cref{prop:solvable} and
$\phi\coloneqq\log\!\big[(c_{1}'-c\,c_{h})^{-1}\big]
+\log\!\big(1+\tfrac{2S}{\delta R}\sqrt{(d+1)/(2\pi)}\big)
=\Theta\!\big(\log\tfrac{d\,L_{1}L_{2}\Delta_{f}}{\epsH\,\delta}\big)$,
\[
  T=\Theta\!\Big(\frac{L_{1}\phi}{\epsH}\Big),\quad
  S=\Theta\!\Big(\frac{\epsH}{L_{1}L_{2}\phi}\Big),\quad
  F=\Theta\!\Big(\frac{\epsH^{3}}{L_{1}^{2}L_{2}^{2}\phi^{3}}\Big).
\]
\end{proposition}
 
\begin{proof}
By \Cref{lem:escapeeig}, $\mup\ge1+\alpha\epsH$ with $\alpha=1/L_{1}$,
and $\log(1+x)\ge x/2$ on $[0,1]$ gives
$\log\mup\ge\tfrac12\alpha\epsH$, so by~\eqref{eq:C5}
$T\le2\phi/(\alpha\epsH)=\Theta(L_{1}\phi/\epsH)$. Then
$S=\Theta(\Lambda/(\sqrt{1+\alpha^{2}}\,L_{2}T))=\Theta(\epsH/(L_{1}L_{2}\phi))$,
since $\Lambda$ and $\sqrt{1+\alpha^{2}}$ are constants in the tolerances, and~\eqref{eq:C1} gives
$F=\Theta\!\big(L_{1}S^{2}/T\big)
=\Theta\!\big(L_{1}\cdot\tfrac{\epsH^{2}}{L_{1}^{2}L_{2}^{2}\phi^{2}}
\cdot\tfrac{\epsH}{L_{1}\phi}\big)
=\Theta(\epsH^{3}/(L_{1}^{2}L_{2}^{2}\phi^{3}))$.
\end{proof}
 
\subsection{Counting the iterations}
 
\begin{proof}[Proof of \Cref{thm:main}]
We split a run into large-gradient steps and escape episodes.
 
\emph{Large-gradient phase.} While $\norm{g_{k}}\ge\epsg$, no
perturbation is triggered and \Cref{lem:descent} gives
$f(x_{k})-f(x_{k+1})\ge\tfrac{C(\nu)}{L_{1}}\norm{g_{k}}^{2}
>\tfrac{C(\nu)}{L_{1}}\epsg^{2}$. Since the total decrease available is
$\Delta_{f}$, the number of such steps is at most
\begin{equation}\label{eq:Ndesc}
  N_{\mathrm{desc}}\le\frac{L_{1}\Delta_{f}}{C(\nu)\,\epsg^{2}}
  =O\!\big(\epsg^{-2}\big).
\end{equation}
 
\emph{Escape phase.} Each perturbation episode that does not certify a
$2$OSP produces, by \Cref{lem:escape}, a net decrease of at least
$(1-\eta)F$. Here \Cref{lem:escape} is applied to the momentum window of each episode
(\Cref{def:window}). Hence the number of episodes is at most
$\bar K\le\Delta_{f}/((1-\eta)F)+1$, and by the scale of $F$ in
\Cref{prop:scales},
\[
  \bar K=\Theta\!\Big(\frac{L_{1}^{2}L_{2}^{2}\Delta_{f}\,\phi^{3}}{\epsH^{3}}\Big)
  =\Otil\!\big(\epsH^{-3}\big).
\]
Each episode costs $T$ iterations, so
\begin{equation}\label{eq:Nesc}
  N_{\mathrm{esc}}\le T\,\bar K
  =\Theta\!\Big(\frac{L_{1}\phi}{\epsH}\Big)
  \cdot\Theta\!\Big(\frac{L_{1}^{2}L_{2}^{2}\Delta_{f}\phi^{3}}{\epsH^{3}}\Big)
  =\Theta\!\Big(\frac{L_{1}^{3}L_{2}^{2}\Delta_{f}\,\phi^{4}}{\epsH^{4}}\Big)
  =\Otil\!\big(\epsH^{-4}\big).
\end{equation}
 
\emph{Union bound.} Index the episodes by $j=1,\dots,\bar K$; episode
$j$ takes a snapshot $\xt^{(j)}$ with $\norm{\grad f(\xt^{(j)})}\le\epsg$
and runs a window of length $T$. Define the bad event
$\mathcal{E}_{j}=\{\lmin(\Hess f(\xt^{(j)}))\le-\epsH
\text{ and the window fails to decrease by }F\}$; that is, $\xt^{(j)}$
was a strict saddle from which the method did not escape. By
\Cref{lem:escape}, $\Prob(\mathcal{E}_{j})\le\delta$ for every $j$.
Fixing a total budget $\delta_{\mathrm{tot}}\in(0,1)$ and setting
$\delta=\delta_{\mathrm{tot}}/\bar K$, the union bound gives
\[
  \Prob\Big(\textstyle\bigcup_{j=1}^{\bar K}\mathcal{E}_{j}\Big)
  \le\sum_{j=1}^{\bar K}\Prob(\mathcal{E}_{j})
  \le\bar K\,\delta=\delta_{\mathrm{tot}}.
\]
With probability at least $1-\delta_{\mathrm{tot}}$, then, no saddle
episode fails. This choice only affects the logarithmic factor: since
$\log(1/\delta)=\log(\bar K/\delta_{\mathrm{tot}})
=\log\bar K+\log(1/\delta_{\mathrm{tot}})$ and
$\log\bar K=\Theta(\log(1/\epsH))=O(\phi)$, it is absorbed into $\phi$,
preserving the scales of \Cref{prop:scales}.
 
Summing~\eqref{eq:Ndesc} and~\eqref{eq:Nesc} yields
$N=O(\epsg^{-2})+\Otil(\epsH^{-4})$, completing the proof.
\end{proof}
 
\begin{table}[t]
\centering
\renewcommand{\arraystretch}{1.35}
\begin{tabular}{@{}lll@{}}
\hline
Symbol & Meaning & Definition / constraint\\
\hline
$\nu$ & restart slack & $0<\nu<\sqrt5-2$\\
$C(\nu)$ & decrease constant & $(1-\nu)-\tfrac12(1+\nu)^{2}=\tfrac12(1-4\nu-\nu^{2})$\\
$\beta$ & momentum (constant) & $\beta\in(0,1]$\\
$\theta$ & warm-up margin & $0<\theta<1$\\
$c$ & residual fraction & $0<c<\min\{1,c_{1}'/c_{h}\}$\\
$\eta$ & perturbation-cost fraction & $0<\eta\le\tfrac12$\\
$\delta$ & per-episode failure prob.\ & $\delta=\delta_{\mathrm{tot}}/\bar K$\\
$\epsg$ & gradient tolerance & $\epsg>0$\\
$\epsH$ & curvature tolerance & $0<\epsH\le L_{1}$\\
$\alpha$ & step size & $1/L_{1}$\\
$H,\lambda_{i}$ & $\Hess f(\xt)$ and its eigenvalues & $\lambda_{\min}\le-\epsH$\\
$B_{i}$ & $\bigl(\begin{smallmatrix}1-\alpha\lambda_{i}&\alpha\beta\\-\lambda_{i}&\beta\end{smallmatrix}\bigr)$ & $\det B_{i}=\beta$\\
$\mu_{\pm}$ & eigenvalues of $\Bmin$ & $\mup\ge1+\alpha\epsH>1>\mum$\\
$\rho_{i}$ & spectral radius of $B_{i}$ & $\max\{\abs{\mu_{1}},\abs{\mu_{2}}\}\le\mup$\\
$C(B)$ & power constant & $2(\norm{B}+\rho)/\abs{\mu_{1}-\mu_{2}}$\\
$c_{h},C_{0}$ & block constants & $c_{h}=C(\Bmin)\le C_{0}=\max_{i}C(B_{i})$\\
$c_{1}$ & energy growth & $(1-\theta)\norm{\vp}/(\mup-\mum)$\\
$c_{1}'$ & position growth & $(1-\theta)\abs{\beta-\mup}/(\mup-\mum)$\\
$\sg$ & residual coefficient & $\sqrt{1+\alpha^{2}}\,L_{2}(S+R)$\\
$S$ & localization radius & $(1+\nu)\sqrt{TF/(L_{1}C(\nu))}$\\
$F$ & escape decrease & $L_{1}C(\nu)\,S^{2}/((1+\nu)^{2}T)$\\
$R$ & perturbation radius & \Cref{prop:solvable}\\
$R_{0},\Rstar$ & separation, threshold & $\Rstar=2S/((c_{1}'-c\,c_{h})\mup^{T}-1)$\\
$T$ & escape window & \eqref{eq:C5}; $\Theta(L_{1}\phi/\epsH)$\\
$\phi$ & log factor & $\Theta(\log(d L_{1}L_{2}\Delta_{f}/(\epsH\delta)))$\\
$\bar K$ & episode count & $\Otil(\epsH^{-3})$\\
\hline
\end{tabular}
\caption{Free parameters, problem data and derived quantities.}
\label{tab:ledger}
\end{table}
 
% ============================================================
\section{Conclusion}
\label{sec:conclusion}
 
We have shown that a perturbed, restarted constant-momentum
method---equivalently, a safeguarded heavy-ball iteration---escapes
strict saddle points and reaches an $(\epsg,\epsH)$-second-order
stationary point in $O(\epsg^{-2})+\Otil(\epsH^{-4})$ iterations,
matching the rate of perturbed gradient descent. The analysis transfers
the perturbation-and-coupling framework of~\cite{pmlr-v70-jin17a,pmlr-v75-jin18a} to an
iteration whose linearised propagator is non-symmetric, via an explicit
spectral study of the two-trajectory difference.
 
The momentum coefficient $\beta$ was held constant throughout, and this
is the main limitation of the present result. The natural next step, and
our motivating goal, is the genuine nonlinear conjugate gradient method,
in which $\beta_{k}=\norm{g_{k}}^{2}/\norm{g_{k-1}}^{2}$
(Fletcher--Reeves~\cite{Fletcher1964FunctionMB}) varies with the iterates
and is singular at stationary points, where the denominator
$\norm{g_{k-1}}^{2}$ vanishes. That singularity is precisely what
obstructs a direct extension of the propagator analysis---the block
$\Bmin$ would itself depend on $k$ through $\beta_{k}$---and treating it
appears to require new ideas. Another interesting future direction would consist in adapting the argument to a line-search
regime, where Wolfe conditions should supply the angle and decrease
guarantees that the fixed step size and restart test furnish here.
 
% ============================================================
\bibliographystyle{plain}
\bibliography{refs}
 
\end{document}